% Identification: losses, optimization, stopping rule, model selection.
% Numbers are measured from the runs recorded in experiments/ledger.md.

\section{Identification}
\label{sec:ident}

\subsection{Objective}
\label{sec:ident-loss}

Let $\Theta$ collect all identified quantities: the residual-source network
weights, the three transport-coefficient scalars, the eight latent
parameters (topology $\beta$, time scale $\tau$, four drive coefficients,
readout gain and offset), and the observation-head weights. Writing
$\phi_\delta(x) = \delta^2\big(\sqrt{1 + (x/\delta)^2} - 1\big)$ for the
pseudo-Huber function, the per-discharge objective is
\begin{equation}
\mathcal{L}(\Theta) =
\underbrace{\frac{\sum_{ij} W_{ij}\,\phi_1\!\big(r_{ij}/s_T\big)}{\sum_{ij} W_{ij}}}_{\text{profile}}
+ \lambda_s\,\overline{\phi_1\!\big(s_\theta/s_s\big)}
+ \lambda_1\,\overline{\lat^{2}}
+ \lambda_2\,\overline{\dot\lat^{2}}
+ \lambda_r\,\overline{\mathrm{BCE}\big(\pH,\ell\big)}
+ \lambda_y\,\overline{\big(\hat\yD - \tilde\yD\big)^{2}},
\label{eq:loss-full}
\end{equation}
with $r_{ij} = \hat T_{\mathrm{e}}(\rho_i,t_j) - T_{\mathrm{e},ij}$ the
profile residual. Each term is an average of order-one quantities, so the
weights are interpretable ratios rather than scale compensators. The fixed
robust scales are $s_T = \SI{100}{eV}$ (comparable to the measurement
scatter of the outer Thomson channels) and $s_s = \SI{1e4}{eV/s}$; both are
frozen constants, not per-discharge statistics.

\paragraph{Which samples are scored.} The profile term is evaluated only
where a measurement exists and survives the quality gates, and only in the
reliable annulus. Its weights $W_{ij}$ are inverse-coverage weights: a
radial channel measured in few time slices receives proportionally more
weight per sample, so densely and sparsely observed radii contribute
comparably instead of the well-covered channels dominating. The regime term
is scored only on samples labelled cleanly L or H; transition windows and
out-of-gate samples are excluded entirely, so ambiguous instants can neither
reward nor penalise the model.

\paragraph{Weights.} $\lambda_s = 0.1$, $\lambda_1 = \lambda_2 = 10^{-3}$,
$\lambda_r = 0.5$, $\lambda_y = 2$. The latent regularisers are deliberately
weak: $\lambda_1$ penalises latent magnitude and $\lambda_2$ its rate, and
both are set an order of magnitude below the terms that carry information,
because a strong smoothness prior would suppress precisely the fast
switching the model is meant to identify. We report in
\S\ref{sec:res-modelcomp} that the topology conclusion is unchanged when
these are varied by a factor of ten.

\subsection{Optimisation and the stopping rule}
\label{sec:ident-opt}

Parameters are updated with AdamW (peak learning rate $2\times10^{-4}$,
cosine decay, five warm-up steps, global-norm gradient clipping at
$10^{3}$, weight decay $3\times10^{-5}$) on mini-batches of eight
discharges drawn uniformly from the training sessions, in double precision.
Any step whose gradient contains a non-finite entry is skipped rather than
applied, which makes a single pathological rollout unable to destroy an
otherwise converged fit; across all runs reported here the skip rate was
zero.

Training stops when the grouped-validation loss has failed to improve by at
least $5\times10^{-4}$ relative for sixteen consecutive evaluations
(evaluations are $50$ steps apart, so the patience is $800$ steps), subject
to a ceiling of $3000$ steps. We record the empirical behaviour of this rule
because it is easy to set badly. Across the nine main runs the best
checkpoint occurred at a median of step $1850$ (range $1250$--$2500$); only
two runs reached the ceiling, and their validation losses ($0.432$, $0.473$)
sit among their siblings rather than below them, so the ceiling is not
truncating a still-improving fit. Measuring the intervals between successive
validation improvements across all runs ($n = 198$) gives a median of $50$
steps and a maximum of $750$; no improvement ever followed a gap larger than
$800$ steps, which is the justification for that patience value.

\paragraph{An incidental diagnostic.} Under the erroneous flux coordinate of
\S\ref{app:eq-rho}, the same rule stopped runs at best checkpoints of step
$150$--$250$: the fit plateaued almost immediately. After the coordinate was
corrected the best checkpoint moved to $\approx 1850$. The location of the
plateau is therefore itself a diagnostic of whether the model has real
structure to fit, and we report it as such.

\subsection{Model selection and calibration}
\label{sec:ident-select}

The validation loss used for checkpoint selection is the full objective
\eqref{eq:loss-full}, including the regime and observation terms, so
selection reflects regime fidelity and not only profile fit. Sessions
assigned to the locked test set take no part in training, checkpoint
selection, calibration, or any modelling decision reported here.

The regime probability requires a separate calibration step, and we state
plainly why. Trained end-to-end, $\pH$ discriminates well but is poorly
scaled: pooled over held-out discharges its expected calibration error is
$0.14$--$0.33$ depending on the configuration, and because the raw
probabilities rarely cross $0.5$, threshold-dependent scores such as $F_1$
collapse to zero even where the ranking is near-perfect. We therefore fit a
two-parameter Platt scaling, $\pH \mapsto \sigma(A\,\eta + B)$ with $\eta$
the regime logit, by minimising binary cross-entropy \emph{on the validation
discharges only}, and apply the fitted $(A,B)$ unchanged elsewhere. This
reduces the expected calibration error to $0.036$--$0.095$. Being a monotone
map, it leaves AUC exactly invariant; only the threshold-dependent and
calibration-sensitive scores change. Calibration is thus \emph{measured},
never assumed, and the uncalibrated values are reported alongside.

\subsection{Cost}
\label{sec:ident-cost}

A single identification run takes $2$--$3$ hours on fourteen CPU cores in
double precision; three runs execute concurrently with pinned thread counts.
A consumer GPU gives no advantage here, because the pipeline is float64
throughout and the rollout is latency- rather than throughput-bound; we
measured a tie at the usable batch size and instability in the compiler at
larger ones. The complete experiment grid reported in
\S\ref{sec:res-modelcomp} is thirteen runs.
